\setcounter{definition}{0}
\setcounter{theorem}{0}
\subsection{Kongruenciák}
\subsubsection{Fogalma}
\begin{definition}
	Adott $n \neq 0$ és $a, b$ egészek esetén
	$a$ kongruens $b$-vel modulo $n$, ha $n \mid a - b$.
\end{definition}
\begin{remark*}
	Jelölése: $a \equiv b \mod n$
\end{remark*}
\subsubsection{Tulajdonságai}
\begin{theorem}
	A kongruencia ekvivalencia reláció.
\end{theorem}
\begin{proof}
	\begin{equation*}
		n \mid a - a = 0 \implies a \equiv a \mod n \implies \text{reflexív} \\
	\end{equation*}
	\begin{multline*}
		a \equiv b \mod n \land b \equiv c \mod n \implies
		n \mid a - b \land n \mid b - c \implies \\
		n \mid (a - b) + (b - c) = a - c \implies
		a \equiv c \mod n \implies \text{tranzitív}
	\end{multline*}
	\begin{multline*}
		a \equiv b \mod n \implies n \mid a - b \implies
		n \mid (-1) \cdot (a - b) = b - a \implies \\
		b \equiv a \mod n \implies \text{szimmetrikus}
	\end{multline*}
\end{proof}
\begin{theorem}
	Legyenek $a, b, c, d, n \in \mathbb{Z}, n \neq 0$ és
	$a \equiv b \mod n$ és $c \equiv d \mod n$. Ekkor
	\begin{itemize}
		\item $a + c \equiv b + d \mod n$
		\item $a \cdot c \equiv b \cdot d \mod n$
	\end{itemize}
\end{theorem}
\begin{proof}
	\begin{multline*}
		a \equiv b \mod n \land c \equiv d \mod n \implies
		n \mid a - b \land n \mid c - d \implies \\
		n \mid (a - b) + (c - d) = (a + c) - (b + d) \implies
		a + c \equiv b + d \mod n
	\end{multline*}
	\begin{multline*}
		a \equiv b \mod n \land c \equiv d \mod n \implies
		n \mid a - b \land n \mid c - d \implies \\
		n \mid (a - b) \cdot c + (c - d) \cdot b =
		a \cdot c - b \cdot d \implies
		a \cdot c \equiv b \cdot d \mod n
	\end{multline*}
\end{proof}
\begin{theorem}
	Legyenek $a, b, c, n \in \mathbb{Z}, n \neq 0$. Ekkor
	\begin{equation*}
		a \cdot b \equiv a \cdot c \mod n \iff b \equiv c \mod \frac{n}{(a, n)}
	\end{equation*}
\end{theorem}
\begin{proof}
	\begin{equation*}
		d = (a, n)
	\end{equation*}
	\begin{multline*}
		a \cdot b \equiv a \cdot c \mod n \implies
		n \mid a \cdot b - a \cdot c = a \cdot (b - c) \implies \\
		\frac{n}{d} \cdot d \mid
		\frac{a}{d} \cdot d \cdot (b - c) \implies
		\exists k \in \mathbb{Z}: k \cdot \frac{n}{d} \cdot d =
		\frac{a}{d} \cdot d \cdot (b - c) \implies \\
		k \cdot \frac{n}{d} = \frac{a}{d} \cdot (b - c) \implies
		\frac{n}{d} \mid \frac{a}{d} \cdot (b - c)
	\end{multline*}
	\begin{equation*}
		\left(\frac{n}{d}, \frac{a}{d}\right) = 1 \implies
		\frac{n}{d} \mid (b - c)
	\end{equation*}
\end{proof}
